\section{Benchmarking} \label{sec:9_benchmarking}
\captionsetup[table]{justification=centering,singlelinecheck=false}
\captionsetup[figure]{justification=centering,singlelinecheck=false}
\newcommand{\thead}[1]{\shortstack[c]{\bfseries #1}}

En esta sección se presenta el análisis de los resultados obtenidos a través del módulo
de \textit{benchmarking}. Se detalla el objetivo de dicho análisis, la metodología seguida para
su ejecución, el entorno experimental utilizado y las comparativas realizadas.

\subsection{Objetivo} \label{sec:9_benchmarking_objective}
El objetivo de este análisis es la comparación del rendimiento de distintos
\textit{pipelines} de procesado de imágenes bajo diversas condiciones de ejecución.
Para ello, se tendrá en cuenta tanto la influencia del hardware como las diferentes configuraciones de procesamiento,
con el fin de evaluar su eficiencia en escenarios de procesamiento en tiempo real.
\noindent En particular, el análisis busca responder a las siguientes preguntas:

\begin{enumerate}[noitemsep]
	\item ¿Qué ganancia aporta la ejecución de \textit{pipelines} de procesado en GPU respecto a CPU?
	\item ¿Cómo varía el rendimiento con la resolución de imagen y con el tipo numérico empleado?
	\item ¿Qué impacto tiene el aumento de la complejidad del \textit{pipeline} de procesado?
	\item ¿Qué diferencias de coste temporal presentan los módulos de \textit{Debayer} implementados?
	\item ¿Cuál es el rendimiento en escenarios de uso realistas?
\end{enumerate}

\subsection{Metodología} \label{sec:9_benchmarking_methodology}
Los resultados analizados proceden de los artefactos JSON generados por el módulo
de \textit{benchmarking}, ejecutado sobre tres equipos
independientes y las siguientes definiciones de \textit{pipelines} de procesado:

\begin{itemize}[noitemsep]
	\item \texttt{standard\_run}: \textit{pipeline} simple de complejidad reducida, consta de las siguientes operaciones:
	      \begin{enumerate}[noitemsep]
		      \item \textit{Debayer5x5()}
		      \item \textit{SigmoidContrast(gain=2.25, cutoff=0.25)}
		      \item \textit{RealToRGB8()}
	      \end{enumerate}
	\item \texttt{complex\_run}: \textit{pipeline} complejo con más operaciones, consta de:
	      \begin{enumerate}[noitemsep]
		      \item \textit{Debayer5x5()}
		      \item \textit{MinMaxNormalization()}
		      \item \textit{GammaAdjustment(c=1.65, gamma=0.8)}
		      \item \textit{UnsharpMasking(kernel\_size=9, sigma=3.0)}
		      \item \textit{RealToRGB8()}

	      \end{enumerate}
	\item \texttt{debayer2x2\_run}, \texttt{debayer3x3\_run}, \texttt{debayer5x5\_run},
	      \texttt{debayerSplit\_run}: perfiles dedicados para comparar las variantes de \textit{Debayer} implementadas.
\end{itemize}

\indent{Las ejecuciones se han llevado a cabo sobre imágenes sintéticas considerando todas las combinaciones posibles de los siguientes factores:}

\begin{itemize}[noitemsep]
	\item \textbf{Resolución de imagen}: \texttt{256x256}, \texttt{512x512}, \texttt{1024x1024} y \texttt{2048x2048}.
	\item \textbf{Dispositivo de ejecución}: \texttt{CPU} y \texttt{GPU}.
	\item \textbf{Tipo numérico de datos}: \texttt{float16}, \texttt{float32} y \texttt{float64}.
\end{itemize}

\newpage

\subsection{Entorno experimental} \label{subsec:9_benchmarking_environment}
Los escenarios de evaluación se han ejecutado sobre tres plataformas hardware
independientes, seleccionadas con el objetivo de cubrir un espectro amplio de
contextos de uso. Cada una de ellas representa un perfil distinto de
despliegue. A continuación se describen las características técnicas de cada equipo y el
escenario de uso que representan.

\subsubsection{Equipo 1. Estación de trabajo de alto rendimiento} \label{subsubsec:equipo1}

El Equipo 1 corresponde a una estación de trabajo de
escritorio, representativa de un entorno de desarrollo, experimentación y
validación de algoritmos. Su especificación de hardware se
detalla en la \autoref{tab:equipo1_specs}.

\begin{table}[h!]
	\centering
	\renewcommand{\arraystretch}{1.25}
	\setlength{\tabcolsep}{10pt}
	\begin{tabularx}{0.85\linewidth}{|>{\hsize=0.55\hsize\centering\arraybackslash}X|>{\hsize=0.85\hsize\centering\arraybackslash}X|>{\hsize=0.6\hsize\centering\arraybackslash}X|}
		\hline
		\rowcolor[gray]{0.90}
		\thead{Componente} & \thead{Modelo}         & \thead{Especificaciones}                                                                \\
		\hline
		\rowcolor{white}
		CPU                 & AMD Ryzen 9 9950X3D     & 16 núcleos / 32 hilos, arquitectura x86\_64, 192 MB caché L3 (64 MB + 128 MB 3D V-Cache) \\
		\hline
		\rowcolor{white}
		GPU                 & NVIDIA GeForce RTX 5090 & Arquitectura Blackwell, 32 GB VRAM                                                       \\
		\hline
		\rowcolor{white}
		Memoria RAM         & DDR5                    & 128 GB                                                                                   \\
		\hline
	\end{tabularx}
	\caption{Especificaciones técnicas del Equipo 1}
	\label{tab:equipo1_specs}
\end{table}

Este equipo se emplea como referencia de rendimiento máximo dentro del estudio. Su combinación de componentes lo sitúa en el
extremo superior del espectro de capacidades disponibles para el desarrollo de algoritmos de visión por computador.
En el contexto del \textit{benchmarking}, permite establecer la cota superior de rendimiento alcanzable para
cada \textit{pipeline} y analizar la escalabilidad de las implementaciones en condiciones donde los recursos
hardware no actúan como factor limitante.

\newpage

\subsubsection{Equipo 2. Sistema embebido de gama alta} \label{subsubsec:equipo2}

El Equipo 2 corresponde a una plataforma de desarrollo NVIDIA Jetson AGX Orin,
un sistema embebido diseñado para cargas de trabajo
de inteligencia artificial y procesado de señales en entornos con restricciones
de energía, espacio y disipación térmica. Sus características se recogen en la
\autoref{tab:equipo2_specs}.

\begin{table}[h!]
	\centering
	\renewcommand{\arraystretch}{1.25}
	\setlength{\tabcolsep}{10pt}
	\begin{tabularx}{0.85\linewidth}{|>{\hsize=0.55\hsize\centering\arraybackslash}X|>{\hsize=0.85\hsize\centering\arraybackslash}X|>{\hsize=0.6\hsize\centering\arraybackslash}X|}
		\hline
		\rowcolor[gray]{0.90}
		\thead{Componente} & \thead{Modelo}         & \thead{Especificaciones}                          \\
		\hline
		\rowcolor{white}
		CPU                 & NVIDIA Carmel (ARMv8.2) & 8 núcleos, arquitectura aarch64, 8 MB L2 + 4 MB L3 \\
		\hline
		\rowcolor{white}
		GPU                 & NVIDIA Volta            & 512 núcleos CUDA + 64 núcleos Tensor               \\
		\hline
		\rowcolor{white}
		Memoria RAM         & LPDDR5                  & 64 GB (compartida CPU/GPU)                         \\
		\hline
	\end{tabularx}
	\caption{Especificaciones técnicas del Equipo 2}
	\label{tab:equipo2_specs}
\end{table}

La arquitectura del AGX Orin presenta características diferenciales respecto a
una estación de trabajo convencional. Opera sobre arquitectura ARM. Como elemento distintivo, 
destaca su modelo de memoria unificada. Los 64 GB de RAM son compartidos entre CPU y GPU, eliminando la
necesidad de transferencias explícitas entre memorias y reduciendo la sobrecarga
asociada al movimiento de datos. No obstante, esta memoria opera a anchos de
banda inferiores a los de una GPU discreta de escritorio.

En el contexto del análisis, este equipo representa un escenario de despliegue en entornos
industriales o de robótica avanzada, donde se requiere capacidad de cómputo significativa en
un formato compacto y con consumo energético contenido. Ejemplos incluyen sistemas de inspección
visual en líneas de producción, vehículos autónomos, robótica móvil y plataformas de
adquisición de datos en campo.

\newpage

\subsubsection{Equipo 3. Sistema embebido de gama media} \label{subsubsec:equipo3}

El Equipo 3 se corresponde con una plataforma NVIDIA Jetson Orin Nano, un módulo
de sistema orientado a aplicaciones de inferencia y procesado de señales con
restricciones más severas de coste y consumo energético. Sus especificaciones se
presentan en la \autoref{tab:equipo3_specs}.

\begin{table}[h!]
	\centering
	\renewcommand{\arraystretch}{1.25}
	\setlength{\tabcolsep}{10pt}
	\begin{tabularx}{0.85\linewidth}{|>{\hsize=0.55\hsize\raggedright\arraybackslash}X|>{\hsize=0.85\hsize\centering\arraybackslash}X|>{\hsize=0.6\hsize\centering\arraybackslash}X|}
		\hline
		\rowcolor[gray]{0.90}
		\thead{Componente} & \thead{Modelo}      & \thead{Especificaciones}             \\
		\hline
		\rowcolor{white}
		CPU                 & ARMv8 (Cortex-A78AE) & 6 núcleos, arquitectura aarch64       \\
		\hline
		\rowcolor{white}
		GPU                 & NVIDIA Ampere        & 1024 núcleos CUDA + 32 núcleos Tensor \\
		\hline
		\rowcolor{white}
		Memoria RAM         & LPDDR5               & 8 GB (compartida CPU/GPU)             \\
		\hline
	\end{tabularx}
	\caption{Especificaciones técnicas del Equipo 3}
	\label{tab:equipo3_specs}
\end{table}

El Orin Nano incorpora una GPU basada en arquitectura NVIDIA Ampere con 1024
núcleos CUDA, el doble que el AGX Orin, si bien con un ancho de banda de memoria
y un consumo energético significativamente menores. Los 8 GB de memoria
compartida son un factor bastante limitante para el procesado de imágenes de gran
resolución. Al igual que el AGX Orin, opera sobre arquitectura ARM.

En el contexto de este análisis, este equipo representa un escenario de despliegue
en dispositivos de bajo coste y consumo mínimo, orientados a aplicaciones de visión
integradas en productos comerciales, sensores inteligentes, cámaras de inspección
portátiles y sistemas \textit{edge} en arquitecturas de \textit{edge computing}.

\newpage

